
\chapter*{序}
\addcontentsline{toc}{chapter}{序}
\markboth{序}{序}

两位作者推出新著《数据科学中的靶向学习》（Targeted Learning in Data Science），值得祝贺。本书将两个令人振奋的领域结合起来，为现有文献增添了一部佳作。它是《靶向学习》（Targeted Learning，van der Laan and Rose 2011）的续篇，建立在前书对机器学习与因果推断所作的广泛而成功的论述之上。靶向学习和数据科学都在迅速发展，边界尚不清晰。为了限定讨论范围，我们将从生物医学数据科学家的角度，说明靶向学习如何帮助统一这一领域的基础。我们尤其关注它如何连接机制建模与基于随机化的推断，从而弥合因果研究基础中的一个重要分野。

具体而言，在\textbf{机制建模（mechanistic modeling）}中，研究者力图准确刻画底层数据生成过程。换句话说，他们对单个观测数据单位的整个分布，或其中相关的部分进行建模。例如，可以明确地为给定处理类型和基线协变量时结局的条件分布建立模型。这正是传统似然建模的核心，也包括常用的贝叶斯方法。机制建模即使考虑因果性，通常也只是作非正式的理解，主要注意力仍放在建模或预测上。历史上，分析以简约性为指导，往往通过对条件均值模型作线性假设来实现。近年来，机器学习方法使我们能够构建更准确、但计算更复杂的模型。因此，有人把机制建模方法分为基于模型的方法和算法方法，例如 Breiman et al. (2001) 的讨论。基于模型的机制建模用相对较少的参数描述数据生成过程。尽管它本身存在模型设定错误的风险，这种方法仍可能具有吸引力，因为在不少情况下，简约模型即使有误，也能对数据生成过程的关键特征作出有用的描述。

当研究者关心个体层面的推断时，基于模型的方法可能受到限制：为追求简约而作出的简化，通常与丰富的数据生成过程不相容，也未必能够准确捕捉总体异质性。相比之下，复杂预测算法往往确实有助于在个体层面开展准确推断，从而支持个体决策。然而，这类算法通常无法简洁描述，也不易直接对有用的总体概括量进行推断。不过，无论采用算法预测还是简约统计模型，研究者都是在为数据生成机制建立模型。

另一方面，许多用于因果推断的统计方法，侧重利用辅助随机化（ancillary randomization），或基于辅助随机化的思想来估计边际效应。研究者通常关心\textbf{平均处理效应（average treatment effect，ATE）}。ATE 是一个很有吸引力的估计目标，因为它概括了一项政策的因果影响，例如用某种药物治疗患者平均会带来多大收益。因果推断研究一直致力于减少这样一类假设：较为强调机制、以模型为基础的方法，在估计 ATE 或相关因果估计目标时所需的假设。在这些方法中，为了避免传统方法在获得因果估计时要求的繁重建模假设，研究者需要付出另一种代价，即在分析中正式纳入随机化方案或处理分配模型。

可以说，机制模型中的算法方法和机器学习方法，以及稳健的边际因果效应估计，都在统计学中引发了变革，也构成了近来学术思潮的重要部分。这两类方法各有长处，却很少彼此交汇。靶向学习最具吸引力的一点，就是在形式上将它们统一起来。对 ATE 开展靶向学习，需要同时估计条件结局均值，即结局回归，以及处理分配分布，即倾向评分；同时又通过\textbf{靶向最小损失估计（targeted minimum loss-based estimation，TMLE）}，实现边际因果效应的稳健估计。通过这种统一，靶向学习构建了一套以底层机制模型为基础的因果推断理论，例如 Judea Pearl 为 van der Laan and Rose (2011) 所作的序言中便讨论了这一点。与以实验观念为基础的因果推断研究者相比，数据科学家往往更习惯从机制出发思考，因此会认同这一底层框架。靶向学习的关键在于：虽然它以估计底层机制模型为基础，但要有效发挥作用，并不要求研究者事先知道或假定某些过于简单的模型或模型类。

靶向学习框架统一了这两个分离的领域，而本书推荐的具体实现还提供了一些近乎额外获得的好处，这与为处理分配概率，即倾向评分建模时的情形类似。具体来说，单是对处理与协变量之间的关系进行建模，就能够提供信息。当然，当目标要求为处理分配建模时，这一点同样适用于靶向学习。类似地，为开展靶向学习而建立多个条件结局模型，往往也会极大加深我们对数据和研究情境的理解。换言之，靶向学习并没有排除基于模型的统计学中那些非常有价值、并非直接针对目标参数的非正式学习。作者明确鼓励以十分广泛的方式开展条件建模，并建议用模型堆叠／超级学习（model stacking/super learning）来融合各个估计。在靶向学习框架中，这种做法打破了机制建模中基于模型的方法与算法方法之间的界限，使两者的区分在实际中几乎不再重要。它还消除了模型选择后推断所带来的顾虑，而这正是较为强调机制的方法面临的关键问题。所有相关模型和算法，都可以按照客观、数据驱动但又事先规定的方式，为堆叠预测模型作出贡献。再结合 TMLE 程序中的靶向更新步骤，模型选择导致错误率有偏的问题基本上得以消除。

另一个好处来自对研究者的明确要求：必须真正指定所关心的边际估计目标，也就是靶向学习的目标。这一要求在数学上带来了优势：TMLE 可以调整机制建模步骤的输出，使之聚焦于目标，从而改善估计。普通极大似然估计等方法则优化整体拟合程度，不专门针对某一目标。靶向学习还有实际益处：它能使讨论更加集中，并推动研究者事先提出假设。打个比方，在靶向学习中，不能先射出一支箭，再围绕它落下的位置画靶心，因为目标必须预先指定。需要强调的是，所关心的效应不必在形式上具有因果含义，也可以只是数据生成机制的任意整体概括量。这一点对于神经影像和神经科学十分重要，也与我们自身的一些研究领域密切相关。在这些领域，科学家往往关心疾病状态或脑形态等“处理”，而这些“处理”在概念上并不能被分配。

说到我们的研究领域，靶向学习将机制建模与机器学习结合起来的特点，对神经科学和神经影像数据科学尤其有吸引力，因为机制性、基于模型的因果方法在这些领域占主导地位。许多最常用的技术都以模型为基础，包括动态因果建模、Granger 因果性和结构方程模型（Friston et al. 2013；Zhang et al. 2015；Chén et al. 2017；Friston et al. 2003；Penny et al. 2004）。神经影像文献经常用“生成建模”（generative modeling）一词，来正面说明某种方法的研究动机。我们可以推测：希望从机制上理解大脑及其疾病的目标，使这个领域更倾向于在观测建模中采用强调机制的方法。\footnote{有趣的是，反过来的过程似乎同样存在：利用统计学思想，为神经组织方式提出假说模型。近年来试图把大脑刻画为本质上具有贝叶斯性质的研究，正是这样的例子（Knill and Pouget 2004）。}然而，因果推断受到的关注要少得多，研究中过分强调基于单个最终模型的统计分析，而作者对此的批评是有道理的。值得注意的是，人们对大脑机制的理解，促成了人工神经网络的发展，并通过深度学习推动了现代人工智能。这些灵活方法把大量传统统计建模方法包含为特例，如今已在数据科学中占据主导地位，在技术相关行业尤其如此。可以设想，将靶向学习用于现有人工智能系统，可以即时进行因果分析，替代正式而耗时的 A/B 检验。

总之，对于希望开展因果分析和边际效应靶向估计的现有数据科学家，本书将起到桥梁作用。它也将帮助因果统计学界理解数据科学中的关键问题及其应用。

\begin{flushright}
Brian Caffo\quad 美国马里兰州巴尔的摩\par
Marco Carone\quad 美国华盛顿州西雅图\par
2017 年 12 月
\end{flushright}
